% GENERATED FILE. Edit canonical lessons, then run npm run generate:books.
% canonical-source-hash: fnv1a64:c994d9595976fbd6
% canonical-lessons: PT-C471-sinceridade, PT-C471-franqueza, PT-C471-iludir, PT-C471-mentiroso, PT-C471-fingir

\chapter{Sinceridade e mentiras}
\label{ch:pt-sinceridade-e-mentiras}
\hlchaptermodality{471}

\begin{chapteropening}
\textbf{By the end of this chapter:} \emph{I can talk about honesty, lies and pretending.}

Complete the last lesson of chapter 471: I can talk about honesty, lies and pretending.
\end{chapteropening}

\section[a sinceridade]{a sinceridade — sincerity, honesty}

\label{lesson:PT-C471-sinceridade}



\begin{quote}
Before the new one: say the Portuguese for modesty, then the Portuguese for to praise.
\end{quote}

\subsection*{You'll want to know: a sinceridade}
\textbf{a sinceridade} — \textquotedblleft{}sincerity, honesty\textquotedblright{}.

\begin{quote}
Agradeço a tua sinceridade.
\end{quote}

\begin{quote}
\emph{I appreciate your honesty.}
\end{quote}

The adjective is \emph{sincero}.

\begin{cousinweb}
\emph{sinceridade} is a learned borrowing from Latin \emph{sinceritās} \textquotedblleft{}purity, soundness\textquotedblright{} (Lewis \& Short), built on \emph{sincērus} \textquotedblleft{}pure, whole\textquotedblright{}.
\end{cousinweb}

\begin{grammarlens}[title={What you've built}]
One more word for telling, explaining and describing at B1.
\end{grammarlens}

\subsection*{Your turn}
\begin{itemize}
\raggedright
  \item \emph{Say it:} \emph{a sinceridade}
  \item \emph{Say it:} \emph{a sinceridade}, once more
  \item \emph{Say it:} a sentence of your own with \emph{a sinceridade}
  \item \emph{Recall:} say the Portuguese for modesty, then the Portuguese for to praise, then say \emph{a sinceridade} again
\end{itemize}

\begin{tcolorbox}[breakable,colback=teal!4,colframe=teal!35!black,title={Before you move on}]
What does a sinceridade mean? (\textquotedblleft{}Sincerity, honesty\textquotedblright{}.) Say it once more.
\end{tcolorbox}
\section[a franqueza]{a franqueza — frankness}

\label{lesson:PT-C471-franqueza}



\begin{quote}
Before the new one: say the Portuguese for to praise, then the Portuguese for sincerity.
\end{quote}

\subsection*{You'll want to know: a franqueza}
\textbf{a franqueza} — \textquotedblleft{}frankness\textquotedblright{}.

\begin{quote}
Com toda a franqueza, não gostei do filme.
\end{quote}

\begin{quote}
\emph{In all frankness, I did not like the film.}
\end{quote}

\emph{Com franqueza} — frankly. The adjective is \emph{franco}.

\begin{cousinweb}
\emph{franqueza} is built on \emph{franco} \textquotedblleft{}frank\textquotedblright{}, which came from French \emph{franc} (EtymDB), the name of the Franks, Latin \emph{Franci} (Lewis \& Short); the sense \textquotedblleft{}free, open\textquotedblright{} is usually explained by the Franks being the free men of their realm.
\end{cousinweb}

\begin{grammarlens}[title={What you've built}]
One more word for telling, explaining and describing at B1.
\end{grammarlens}

\subsection*{Your turn}
\begin{itemize}
\raggedright
  \item \emph{Say it:} \emph{a franqueza}
  \item \emph{Say it:} \emph{a franqueza}, once more
  \item \emph{Say it:} a sentence of your own with \emph{a franqueza}
  \item \emph{Recall:} say the Portuguese for to praise, then the Portuguese for sincerity, then say \emph{a franqueza} again
\end{itemize}

\begin{tcolorbox}[breakable,colback=teal!4,colframe=teal!35!black,title={Before you move on}]
What does a franqueza mean? (\textquotedblleft{}Frankness\textquotedblright{}.) Say it once more.
\end{tcolorbox}
\section[iludir]{iludir — to deceive, to fool}

\label{lesson:PT-C471-iludir}



\begin{quote}
Before the new one: say the Portuguese for sincerity, then the Portuguese for frankness.
\end{quote}

\subsection*{You'll want to know: iludir}
\textbf{iludir} — \textquotedblleft{}to deceive, to fool\textquotedblright{}.

\begin{quote}
As fotografias do hotel iludiram os turistas.
\end{quote}

\begin{quote}
\emph{The hotel photos fooled the tourists.}
\end{quote}

\emph{Iludir-se} — to fool yourself.

\begin{cousinweb}
\emph{iludir} is a learned borrowing from Latin \emph{illūdere} \textquotedblleft{}to mock, to make sport of\textquotedblright{} (Lewis \& Short), built on \emph{lūdere} \textquotedblleft{}to play\textquotedblright{}.
\end{cousinweb}

\begin{grammarlens}[title={What you've built}]
One more word for telling, explaining and describing at B1.
\end{grammarlens}

\subsection*{Your turn}
\begin{itemize}
\raggedright
  \item \emph{Say it:} \emph{iludir}
  \item \emph{Say it:} \emph{iludir}, once more
  \item \emph{Say it:} a sentence of your own with \emph{iludir}
  \item \emph{Recall:} say the Portuguese for sincerity, then the Portuguese for frankness, then say \emph{iludir} again
\end{itemize}

\begin{tcolorbox}[breakable,colback=teal!4,colframe=teal!35!black,title={Before you move on}]
What does iludir mean? (\textquotedblleft{}To deceive, to fool\textquotedblright{}.) Say it once more.
\end{tcolorbox}
\section[mentiroso]{mentiroso — liar; lying}

\label{lesson:PT-C471-mentiroso}



\begin{quote}
Before the new one: say the Portuguese for frankness, then the Portuguese for to deceive.
\end{quote}

\subsection*{You'll want to know: mentiroso}
\textbf{mentiroso} — \textquotedblleft{}liar; lying\textquotedblright{}.

\begin{quote}
Ninguém acredita nele, porque é um mentiroso.
\end{quote}

\begin{quote}
\emph{Nobody believes him, because he is a liar.}
\end{quote}

Built on \emph{mentir} \textquotedblleft{}to lie\textquotedblright{}; also an adjective: \emph{uma resposta mentirosa}.

\begin{cousinweb}
\emph{mentiroso} is inherited from Old Portuguese \emph{mentiroso}, built on \emph{mentir} \textquotedblleft{}to lie\textquotedblright{} (EtymDB), from Latin \emph{mentīrī}.
\end{cousinweb}

\begin{grammarlens}[title={What you've built}]
One more word for telling, explaining and describing at B1.
\end{grammarlens}

\subsection*{Your turn}
\begin{itemize}
\raggedright
  \item \emph{Say it:} \emph{mentiroso}
  \item \emph{Say it:} \emph{mentiroso}, once more
  \item \emph{Say it:} a sentence of your own with \emph{mentiroso}
  \item \emph{Recall:} say the Portuguese for frankness, then the Portuguese for to deceive, then say \emph{mentiroso} again
\end{itemize}

\begin{tcolorbox}[breakable,colback=teal!4,colframe=teal!35!black,title={Before you move on}]
What does mentiroso mean? (\textquotedblleft{}Liar; lying\textquotedblright{}.) Say it once more.
\end{tcolorbox}
\section[fingir]{fingir — to pretend}

\label{lesson:PT-C471-fingir}



\begin{quote}
Before the new one: say the Portuguese for to deceive, then the Portuguese for liar.
\end{quote}

\subsection*{You'll want to know: fingir}
\textbf{fingir} — \textquotedblleft{}to pretend\textquotedblright{}.

\begin{quote}
O gato fingiu que estava a dormir.
\end{quote}

\begin{quote}
\emph{The cat pretended it was asleep.}
\end{quote}

\emph{Fingir que} + verb — to pretend that. It has nothing to do with fingers.

\begin{cousinweb}
\emph{fingir} is a learned borrowing from Latin \emph{fingere} \textquotedblleft{}to form, to shape, to devise, to feign\textquotedblright{} (EtymDB).
\end{cousinweb}

\begin{grammarlens}[title={What you've built}]
One more word for telling, explaining and describing at B1.
\end{grammarlens}

\subsection*{Your turn}
\begin{itemize}
\raggedright
  \item \emph{Say it:} \emph{fingir}
  \item \emph{Say it:} \emph{fingir}, once more
  \item \emph{Say it:} a sentence of your own with \emph{fingir}
  \item \emph{Recall:} say the Portuguese for to deceive, then the Portuguese for liar, then say \emph{fingir} again
\end{itemize}

\begin{tcolorbox}[breakable,colback=teal!4,colframe=teal!35!black,title={Before you move on}]
What does fingir mean? (\textquotedblleft{}To pretend\textquotedblright{}.) Say it once more.
\end{tcolorbox}
